\chapter{Benvenuti al Secondo Anno: Metodologia e OOP}
\label{chap:intro-secondo-anno}

\begin{quote}
\itshape
\small
``Scrivere codice che un computer comprende è facile. Scrivere codice che gli esseri umani sanno leggere, mantenere ed estendere nel tempo è la vera arte della programmazione.''
\end{quote}
\vspace{1em}

\noindent\textbf{Corso:} Programmazione 2 -- Laboratorio di Sviluppo Software\\
\textbf{Obiettivo del secondo anno:} Passare dalla programmazione procedurale di base alla progettazione modulare e orientata agli oggetti, acquisendo standard professionali di sviluppo.

\vspace{1.5em}

% ------------------------------------------------------------------------------
\section{Dal Pensiero Algoritmico allo Sviluppo Software}
% ------------------------------------------------------------------------------

Nel corso della prima annualità abbiamo esplorato i mattoni fondamentali dell'informatica: la sequenza logica, le condizioni (\texttt{if/else}), i cicli (\texttt{while}, \texttt{for}) e le prime funzioni.

Nel secondo anno compiamo un salto di qualità decisivo: non ci limiteremo a scrivere script che ``funzionano'', ma impareremo a progettare \textbf{architetture software robuste}, manutenibili e modulari.

\begin{concetto}[L'Approccio del Secondo Anno]
Un programma professionale non è un monolite disordinato di centinaia di righe, ma un sistema ordinato di componenti indipendenti che collaborano tra loro scambiandosi messaggi e dati validati.
\end{concetto}

% ------------------------------------------------------------------------------
\section{I Pilastri del Programma Annuale}
% ------------------------------------------------------------------------------

Il percorso di Programmazione 2 si articolerà attorno a quattro direttrici principali:

\begin{enumerate}
  \item \textbf{Programmazione Orientata agli Oggetti (OOP):}
  \begin{itemize}
    \item Concetto di Classe vs Istanza (Oggetto).
    \item Incapsulamento: attributi, metodi e protezione dello stato interno.
    \item Ereditarietà e Polimorfismo: riutilizzo intelligente del codice.
  \end{itemize}

  \item \textbf{Strutture Dati Dinamiche e Collezioni:}
  \begin{itemize}
    \item Liste avanzate, dizionari (tabelle chiave-valore), insiemi e tuple.
    \item Algoritmi di ricerca e ordinamento a confronto.
  \end{itemize}

  \item \textbf{Persistenza dei Dati e Gestione degli Errori:}
  \begin{itemize}
    \item Lettura e scrittura di file (testo, CSV, JSON).
    \item Gestione professionale delle eccezioni (\texttt{try / catch / finally}) per prevenire crash imprevisti.
  \end{itemize}

  \item \textbf{Progetto Pratico di Fine Anno:}
  \begin{itemize}
    \item Realizzazione in team di un applicativo completo, con interfaccia utente, gestione dello stato e salvataggio dei dati.
  \end{itemize}
\end{enumerate}

% ------------------------------------------------------------------------------
\section{Standard di Laboratorio e Buone Pratiche}
% ------------------------------------------------------------------------------

Durante tutte le sessioni di laboratorio adotteremo rigorosi standard professionali:

\begin{consiglio}[Le Regole d'Oro del Laboratorio]
\begin{enumerate}
  \item \textbf{Nomi parlanti:} I nomi di variabili e funzioni devono spiegare \emph{cosa fanno}, non essere sigle oscure come \texttt{a}, \texttt{x1}, \texttt{temp2}.
  \item \textbf{Una funzione, un solo compito:} Ogni funzione deve avere una responsabilità chiara e circoscritta.
  \item \textbf{Controllo di Versione (Git):} Ogni sessione di lavoro prevede il commit delle modifiche con messaggi chiari e significativi.
\end{enumerate}
\end{consiglio}

% ------------------------------------------------------------------------------
\section{Esercizio di Riscaldamento: Verso gli Oggetti}
% ------------------------------------------------------------------------------

Per riscaldare la logica prima di addentrarci nella OOP formale, analizziamo come rappresentare un'entità del mondo reale.

\begin{esercizio}[Modellare uno Studente]
Immagina di dover gestire i dati di una classe di studenti.
\begin{enumerate}
  \item Quali caratteristiche (dati) identificano univocamente ogni studente?
  \item Quali azioni (operazioni) possono essere compiute su uno studente (es. aggiungere un voto, calcolare la media)?
\end{enumerate}
\end{esercizio}

\begin{soluzione}[Dalla Scheda al Modello Dati]
Nel paradigma a oggetti separeremo nettamente:
\begin{itemize}
  \item \textbf{Stato (Attributi):} \texttt{nome}, \texttt{cognome}, \texttt{matricola}, \texttt{elenco\_voti}.
  \item \textbf{Comportamento (Metodi):} \texttt{aggiungi\_voto(voto)}, \texttt{calcola\_media()}, \texttt{is\_promosso()}.
\end{itemize}
Nel prossimo capitolo tradurremo questo modello nella nostra prima classe software!
\end{soluzione}
